\uline{integer} @@M0@@ \uline{such that, for every} @@M1@@, \uline{there is a section} @@M2@@ \uline{that is a basis of} @@M3@@ \uline{for every} @@M4@@.

There is an integer @@M5@@ and a section @@M6@@ such that, for every @@M7@@, @@M8@@, and such that @@M9@@ is an affine open subset (EGA II 4.5.4). On @@M10@@, there is a section @@M11@@ satisfying the conditions of Lemma 6.7; moreover, there is an integer @@M12@@ such that the section @@M13@@ extends to the whole of @@M14@@. Put @@M15@@; then @@M16@@; taking the tensor product of the section of @@M17@@ obtained above and @@M18@@ therefore gives a section of @@M19@@ for every @@M20@@. Since, for every @@M21@@, @@M22@@, and the section @@M23@@ is a basis of @@M24@@ for every @@M25@@, the sections of @@M26@@ thus constructed for every @@M27@@ form a basis of @@M28@@ for every @@M29@@.

\sgaStableTarget{sga6.e06.thm.6-9}\textbf{Theorem 6.9.} \uline{Let} @@M30@@ \uline{be a noetherian scheme of dimension} @@M31@@, \uline{having an ample sheaf} @@M32@@. \uline{Then} @@M33@@, \uline{and a fortiori} @@M34@@.

(This last relation is due to J.P. Serre, in the general case considered here).

We shall first show that if @@M35@@ are two invertible sheaves on @@M36@@, with classes @@M37@@, then @@M38@@.

Let @@M39@@ be a closed subset of @@M40@@, @@M41@@ its irreducible components, and for every @@M42@@, let @@M43@@ be a point of @@M44@@. By Lemma 6.8, there are integers @@M45@@ such that, for any integers @@M46@@, there are homomorphisms from @@M47@@ to @@M48@@, @@M49@@ that are isomorphisms at the points @@M50@@; we therefore obtain morphisms:
@@M51@@
